% Part: turing-machines
% Chapter: machines-computations
% Section: unary-numbers

\documentclass[../../../include/open-logic-section]{subfiles}

\begin{document}

\olfileid{tur}{mac}{una}
\olsection{د عددونو يوګونې څرګندونه}

\begin{explain}
ټيورينګ ماشينونه پر خپله پټه د ليکل شوو نښو پر لړيو کار کوي۔ د
ټيورينګ ماشين د الفبا له مخې، د نښو دا لړۍ بېلابېل ننوتونه او وتونه
څرګندولے شي۔ البته، هغه ټيورينګ ماشينونه ځانګړې پاملرنه غواړي چې
\emph{حسابي} تابعې، يعنې د طبيعي عددونو تابعې، محاسبه کوي۔ د مثبتو
صحيح عددونو د څرګندولو يوه ساده لار دا ده چې د يوې نښې~$\TMstroke$
د لړيو په توګه کوډ شي۔ که $n \in \Nat$ وي، نو $\TMstroke^n$ دې هغه
وخت تشه لړۍ وي چې $n = 0$ وي؛ او که نۀ وي، هغه لړۍ دې وي چې کره
$n$ شمېر~$\TMstroke$ نښې لري۔
\end{explain}

\begin{defn}[محاسبه]
ټيورينګ ماشين~$M$ هله او يوازې هله تابع $f\colon \Nat^k \to \Nat$
\emph{محاسبه کوي} چې~$M$ پر دې ننوت ودرېږي:
\[
\TMstroke^{n_1} \TMblank \TMstroke^{n_2} \TMblank \dots \TMblank \TMstroke^{n_k}
\]
او وت ئې $\TMstroke^{f(n_1, \dots, n_k)}$ وي۔
\end{defn}

\begin{prob}
  دا تعريف کړئ چې ټيورينګ ماشين~$M$ کله تابع
  $f\colon \Nat^k \to \Nat^m$ محاسبه کوي۔
\end{prob}

\begin{ex}\ollabel{ex:adder}
\emph{جمع:}
راځئ داسې ماشين جوړ کړو چې تابع $f(n,m) = n + m$ محاسبه کړي۔ د دې
دپاره داسې ماشين پکار دے چې پر پټه د~$n$ او~$m$ اوږدوالي لرونکو د
$\TMstroke$ دوو غونډونو سره پيل وکړي، او د~$n+m$ شمېر~$\TMstroke$
ګانو په يوه غونډ ودرېږي۔ د~$\TMstroke$ د ننوت دواړه غونډونه په يوه
$\TMblank$ سره بېلېږي؛ نو يوه لار دا ده چې په هغه مربع کښې يوه کرښه
وليکو چې~$\TMblank$ لري، او وروستۍ~$\TMstroke$ پاکه کړو۔
\begin{figure}\[
\begin{tikzpicture}[->,>=stealth',shorten >=1pt,auto,node distance=2.8cm,
                    semithick]
  \tikzstyle{every state}=[fill=none,draw=black,text=black]

  \node[initial,state]         (A)              {$q_0$};
  \node[state]         (B) [right of=A] {$q_1$};
  \node[state]         (C) [right of=B] {$q_2$};

  \path (A) edge node {\TMtrans{\TMblank}{\TMstroke}{\TMstay}} (B)
           edge [loop above] node {\TMtrans{\TMstroke}{\TMstroke}{\TMright}} (B)
        (B) edge [loop above] node {\TMtrans{\TMstroke}{\TMstroke}{\TMright}} (B)
            edge node {\TMtrans{\TMblank}{\TMblank}{\TMleft}} (C)
        (C) edge [loop above] node {\TMtrans{\TMstroke}{\TMblank}{\TMstay}} (C);
\end{tikzpicture}
\]\caption{هغه ماشين چې $f(x,y) = x+y$ محاسبه کوي}
\ollabel{fig:adder}
\end{figure}
\end{ex}

\begin{prob}
په~\olref[tur][mac][una]{ex:adder} کښې د ماشين تشکيلونه د
ننوت~$\tuple{3,2}$ دپاره يو په بل پسې وڅارئ۔ که ماشين $0+0$ محاسبه
کړي، څه پېښېږي؟
\end{prob}

\begin{explain}
په~\olref[rep]{ex:doubler} کښې مو د داسې ټيورينګ ماشين بېلګه ورکړه
چې د~$\TMstroke$ نښو يوه لړۍ د ننوت په توګه اخلي او پر پټه د هغې د
دوه‌چنده شمېر~$\TMstroke$ نښو په لړۍ درېږي---يعنې دوه‌چنده کوونکے
ماشين۔ خو څرنګه چې په وت کښې د~$\TMstroke$ نښو د دوه‌چنده غونډ کيڼ
لوري ته~$\TMblank$ نښې شته، نو دا په حقيقت کښې تابع $f(x) = 2x$ نۀ
محاسبه کوي، که څه هم ښايي تاسو داسې ګڼلې وي۔ موږ به د دې د سمولو
دوه لارې وښيو۔
\end{explain}

\begin{ex}
په~\olref{fig:doubler-disc} کښې ماشين تابع $f(x) = 2x$ محاسبه کوي۔
دا ماشين د~\olref[rep]{ex:doubler} د ماشين په شان د ننوت د هرې
$\TMstroke$ په بدل کښې ننوت نۀ پاکوي او په لرې ښي لوري کښې دوه
$\TMstroke$ نښې نۀ ليکي، بلکې د ننوت د هرې~$\TMstroke$ دپاره ښي
لوري ته يوه~$\TMstroke$ ور زياتوي۔ دا بايد په ياد وساتي چې ننوت چېرته
ختمېږي؛ نو د ننوت او ور زياتو شوو کرښو تر منځ يوه~$\TMblank$ پرېږدي
او په پای کښې ئې په~$\TMstroke$ ډکوي۔ موږ بايد دا هم ``په ياد'' ولرو
چې په ننوت کښې چېرته يو؛ نو د ننوت په غونډ کښې يوه~$\TMstroke$ په
لنډمهاله توګه په~$\TMblank$ بدلوو۔
  \begin{figure}
\[
\begin{tikzpicture}[->,>=stealth',shorten >=1pt,auto,node distance=2.8cm,
                    semithick]
  \tikzstyle{every state}=[fill=none,draw=black,text=black]

  \node[initial,state] (0)              {$q_0$};
  \node[state]         (1) [above of=0] {$q_1$};
  \node[state]         (2) [above of=1] {$q_2$};
  \node[state]         (3) [right of=2] {$q_3$};
  \node[state]         (4) [below of=3] {$q_4$};
  \node[state]         (5) [below of=4] {$q_5$};
  \node[state]         (6) [above right of=4] {$q_6$};
  \node[state]         (7) [below of=6] {$q_7$};
  \node[state]         (8) [below of=7] {$q_8$};

  \path (0) edge node {\TMtrans{\TMstroke}{\TMblank}{\TMright}} (1)
%    (0) edge node {\TMtrans{\TMblank}{\TMblank}{\TMstay}} (h)
    (1) edge [loop left] node {\TMtrans{\TMstroke}{\TMstroke}{\TMright}} (1)
      edge node {\TMtrans{\TMblank}{\TMblank}{\TMright}} (2)
    (2) edge [loop above] node {\TMtrans{\TMstroke}{\TMstroke}{\TMright}} (2)
        edge node {\TMtrans{\TMblank}{\TMstroke}{\TMleft}} (3)
    (3) edge [loop above] node {\TMtrans{\TMstroke}{\TMstroke}{\TMleft}} (3)
        edge node[left] {\TMtrans{\TMblank}{\TMblank}{\TMleft}} (4)
    (4) edge        node[left] {\TMtrans{\TMstroke}{\TMstroke}{\TMleft}} (5)
    (5) edge [loop below] node {\TMtrans{\TMstroke}{\TMstroke}{\TMleft}} (5)
        edge              node {\TMtrans{\TMblank}{\TMstroke}{\TMright}} (0)
    (4) edge      node[sloped] {\TMtrans{\TMblank}{\TMstroke}{\TMright}} (6)
    (6) edge              node {\TMtrans{\TMblank}{\TMstroke}{\TMright}} (7)
    (7) edge [loop right] node {\TMtrans{\TMstroke}{\TMstroke}{\TMright}} (7)
        edge              node {\TMtrans{\TMblank}{\TMblank}{\TMleft}} (8)
    (8) edge [loop right] node {\TMtrans{\TMstroke}{\TMblank}{\TMstay}} (8);
    \end{tikzpicture}
\]
\caption{هغه ماشين چې $f(x) = 2x$ محاسبه کوي}
\ollabel{fig:doubler-disc}
\end{figure}
\end{ex}

\begin{ex}\ollabel{ex:mover}
د تابع $f(x) = 2x$ د محاسبې دويمه شونتيا دا ده چې اصلي دوه‌چنده
کوونکے ماشين وساتو، خو په پای کښې داسې حالتونه او لارښوونې ور زياتې
کړو چې د کرښو دوه‌چنده غونډ د پټې تر ټولو کيڼ لوري ته يوسي۔ په
\olref{fig:mover} کښې ماشين همدا وروستۍ برخه ترسره کوي: که دا پر داسې
پټه پيل شي چې لومړے د~$\TMblank$ نښو يو غونډ او بيا د~$\TMstroke$
نښو يو غونډ لري (او سر د~$\TMblank$ نښو د غونډ په هر ځاے کښې وي)،
نو~$\TMstroke$ نښې يوه يوه پاکوي او د پټې په پيل کښې ئې ليکي۔ د دې
د معلومولو دپاره چې کار کله بشپړ شو، ماشين لومړے د~$\TMstroke$ نښو
د غونډ پای په~$\TMendtape$ نښه نښانوي، او دا نښه په پای کښې پاکېږي۔
موږ د حالتونو شمېرل له~$q_6$ څخه پيل کړي، څو دوي دوه‌چنده کوونکي
ماشين ته ور زياتېدے شي۔ تاسو يوازې يوې اضافي لارښوونې ته اړتيا لرئ:
$\delta(q_0, \TMblank) = \tuple{q_6,\TMblank,\TMstay}$؛ يعنې له~$q_0$
څخه~$q_6$ ته يو داسې غشی چې نښه ئې~$\TMtrans{\TMblank}{\TMblank}{\TMstay}$
وي۔ (يوه نرۍ ستونزه شته: لاس ته راغلے ماشين د ننوت~$x=0$ دپاره کار
نۀ کوي۔ موږ دا د تمرين په توګه پرېږدو۔)
\begin{figure}
  \[
  \begin{tikzpicture}[->,>=stealth',shorten >=1pt,auto,node distance=2.8cm,
                      semithick]
    \tikzstyle{every state}=[fill=none,draw=black,text=black]
    \node[initial,state] (6)              {$q_6$};
    \node[state]         (7) [right of=6] {$q_7$};
    \node[state]         (8) [right of=7] {$q_8$};
    \node[state]         (9) [below of=8] {$q_9$};
    \node[state]         (10) [left of=9] {$q_{10}$};
    \node[state]         (11) [left of=10]  {$q_{11}$};
    \node[state]         (12) [below of=11]  {$q_{12}$};
    \node[state]         (13) [right of=12] {$q_{13}$};
    \node[state]         (14) [right of=13] {$q_{14}$};
    \path
    (6)  edge node {\TMtrans{\TMblank}{\TMblank}{\TMright}} (7)
    (7)  edge [loop above] node {\TMtrans{\TMblank}{\TMblank}{\TMright}} (7)
         edge node {\TMtrans{\TMstroke}{\TMstroke}{\TMright}} (8)
    (8)  edge [loop above] node {\TMtrans{\TMstroke}{\TMstroke}{\TMright}} (8)
         edge node {\TMtrans{\TMblank}{\TMendtape}{\TMleft}} (9)
    (9)  edge [loop right] node {\TMtrans{\TMstroke}{\TMstroke}{\TMleft}} (9)
         edge node {\TMtrans{\TMblank}{\TMblank}{\TMright}} (10)
    (10) edge node[above] {\TMtrans{\TMstroke}{\TMblank}{\TMleft}} (11)
    (11) edge [loop above] node {\TMtrans{\TMblank}{\TMblank}{\TMleft}} (11)
         edge node[left] {\begin{tabular}{@{}l@{}}
          \TMtrans{\TMendtape}{\TMendtape}{\TMright}\\
          \TMtrans{\TMstroke}{\TMstroke}{\TMright}
         \end{tabular}} (12)
    (12) edge node[] {\TMtrans{\TMblank}{\TMstroke}{\TMright}} (13)
    (13) edge [loop below] node {\TMtrans{\TMblank}{\TMblank}{\TMright}} (13)
         edge node[sloped] {\TMtrans{\TMstroke}{\TMblank}{\TMleft}} (11)
    (13) edge node {\TMtrans{\TMendtape}{\TMblank}{\TMstay}} (14);
  \end{tikzpicture}
  \]
  \caption{د~$\TMstroke$ نښو د يو غونډ کيڼ لوري ته وړل}
  \ollabel{fig:mover}
\end{figure}
\end{ex}

\begin{prob}
په~\olref[tur][mac][una]{ex:mover} کښې مو داسې ماشين بيان کړ چې د
\olref[tur][mac][rep]{fig:doubler} د دوه‌چنده کوونکي ماشين او د
\olref[tur][mac][una]{fig:mover} د لېږدوونکي ماشين له يوځاے کولو څخه
جوړ دے۔ که دا يوځاے شوے ماشين پر ننوت~$x=0$، يعنې پر تشه پټه پيل
کړئ، څه پېښېږي؟ ماشين به څنګه سم کړئ چې په دې صورت کښې د~$2x=0$ له
وت سره ودرېږي؟ (تاسو بايد دا د يوه حالت او يوه حالت بدلون په ور
زياتولو وکړے شئ۔)
\textbf{د ژباړې د نسخې سمون (OLCMP-046):} سرچينې دلته د هغه بل
دوه‌چنده کوونکي شکل حواله ورکړې وه چې لا پخپله بشپړ وت جوړوي او
حالتونه ئې له لېږدوونکي سره ټکر کوي؛ د مخکښېني بيان او د حالتونو له
شمېرنې سره سم د اصلي شپږ-حالته دوه‌چنده کوونکي حواله راوړل شوه۔
\end{prob}

\begin{prob}
\emph{تفريق:} داسې ټيورينګ ماشين طرح کړئ چې د~$n$ او~$m$ اوږدوالي
لرونکې د کرښو دوه ناتشې رشتې د ننوت په توګه ورکړل شي، چې $n > m$،
نو تابع $f(n,m) = n - m$ محاسبه کړي۔
\end{prob}

\begin{prob}
\emph{مساوات:} داسې ټيورينګ ماشين طرح کړئ چې لاندې تابع محاسبه کړي:
\[
\fn{equality}(n,m) =
\begin{cases}
  \text{1} & \text{که~$n = m$ وي} \\
  \text{0} & \text{که~$n \neq m$ وي}
\end{cases}
\]
چې پکښې~$n$ او~$m \in \PosInt$ دي۔
\end{prob}

\begin{prob}
د تابع $\min(x,y)$ د محاسبې دپاره يو ټيورينګ ماشين طرح کړئ، چې
$x$ او~$y$ داسې مثبت صحيح عددونه وي چې پر پټه د~$\TMstroke$ نښو په
رشتو څرګند شوي او په يوه~$\TMblank$ بېل شوي وي۔ تاسو د ماشين په
الفبا کښې اضافي نښې کارولے شئ۔

تابع~$\min$ له خپلو ارزښتمنو څخه تر ټولو کوچنی ارزښت ټاکي؛ نو
$\min(3,5)=3$، $\min(20,16)=16$، او $\min(4,4)=4$، او داسې نور۔
\end{prob}

\begin{defn}
  ټيورينګ ماشين~$M$ هله او يوازې هله جزوي تابع $f\colon \Nat^k
  \pto \Nat$ محاسبه کوي چې
  \begin{enumerate}
    \item که $f(n_1, \dots, n_k) = m$ وي، نو~$M$ پر ننوت
    $\TMstroke^{n_1}\concat\TMblank\concat
    \dots \concat\TMblank\concat\TMstroke^{n_k}$ د وت~$\TMstroke^{m}$ سره درېږي۔
    \item که $f(n_1, \dots, n_k)$ تعريف شوې نۀ وي، نو~$M$ بېخي نۀ
    درېږي، يا داسې درېږي چې معتبر وت ورته نۀ ټاکل کېږي، يا له داسې
    وت سره درېږي چې د~$\TMstroke$ نښو يوازينی غونډ نۀ وي۔
  \end{enumerate}
\textbf{د ژباړې د نسخې سمون (OLCMP-061):} د سرچينې دويم شرط دا
فرضوي چې هره درېدلې چلونه يو وت لري۔ د مخکښېني مشروط وت له تعريف
سره د سمون دپاره، هغه درېدنه هم دلته شامله شوه چې معتبر وت ورته نۀ
ټاکل کېږي۔
\end{defn}

\end{document}
